% \pdfstrcmp{A}{B} -- both arguments fully expanded, then compared as strings:
% -1, 0 or 1, by character code. expl3 does not merely prefer it, it binds
% \__str_if_eq:nn straight to it with no fallback, so every string comparison
% in the LaTeX kernel needs it.
\input plain
\catcode`\@=11
\def\M#1{\message{<#1>}}
\M{eq=\pdfstrcmp{abc}{abc}}
\M{lt=\pdfstrcmp{abc}{abd}}
\M{gt=\pdfstrcmp{abd}{abc}}
\M{prefix=\pdfstrcmp{ab}{abc}}
\M{empty=\pdfstrcmp{}{}}
\M{emptyb=\pdfstrcmp{}{a}}
% The arguments are expanded, as in an \edef.
\def\foo{XY}
\M{expands=\pdfstrcmp{\foo}{XY}}
% Control sequences compare by their printed form.
\M{cs=\pdfstrcmp{\relax}{\relax}}
\M{csdiff=\pdfstrcmp{\relax}{\hbox}}
% By character code, so digits sort as text and case matters.
\M{digits=\pdfstrcmp{10}{9}}
\M{case=\pdfstrcmp{A}{a}}
% And it is expandable: usable inside \edef and in a number scan.
\edef\r{\pdfstrcmp{a}{b}}\M{edef=\meaning\r}
\ifnum\pdfstrcmp{a}{a}=0 \M{innum=yes}\else\M{innum=no}\fi
\bye
